\documentclass[11pt]{article}
%%%%%%%%%%%%%%%%%%%%%%%%%%%%%%%%%%%%%%%%%%%%%%%%%%%%%%%%%%%%%%%%%%%%%%%%%%%%%%%%%%%%
% Internal revision TO-DO checklist -- companion to response.tex.
%
% The items are ordered by WHERE the change is made in the manuscript (front to
% back), so this doubles as a revision work-plan: go through the paper top to
% bottom and tick each box.
%   Tags:  [R1.x]/[R2.x] = which referee comment it answers;
%          (minor)/(clear) = quick, self-contained edits;
%          (substantial)   = needs new text, examples, or analysis.
%
% NOTE: compile with the working directory at the PROJECT ROOT (as Overleaf does),
% so the manuscript cross-references (\ref, \pageref) resolve from the committed
% label snapshot (Statistics_and_computing_submission/ustats_revision_labels.aux).
% Each \ref/\pageref shows the CURRENT manuscript number/page.
%%%%%%%%%%%%%%%%%%%%%%%%%%%%%%%%%%%%%%%%%%%%%%%%%%%%%%%%%%%%%%%%%%%%%%%%%%%%%%%%%%%%
\usepackage[T1]{fontenc}
\usepackage{lmodern}
\usepackage{amsmath,amssymb}
\usepackage{enumitem}
\usepackage[svgnames]{xcolor}
\usepackage[colorlinks=true, linkcolor=teal, citecolor=blue, urlcolor=blue]{hyperref}

% manuscript cross-references (same robust setup as response.tex)
\usepackage{xr-hyper}
\IfFileExists{ustats_revision.aux}{%
  \externaldocument{ustats_revision}%
}{\IfFileExists{Statistics_and_computing_submission/ustats_revision.aux}{%
  \externaldocument{Statistics_and_computing_submission/ustats_revision}%
}{%
  \externaldocument{Statistics_and_computing_submission/ustats_revision_labels}%
}}

\newcommand{\bx}{$\square$\ }   % an open checkbox; change to $\boxtimes$ when done
\newcommand{\at}[1]{\textbf{#1}} % a manuscript location

\title{Revision To-Do Checklist\\[-0.2em]
\large (companion to \texttt{response.tex}, ordered by manuscript location)}
\date{\today}
\author{}

\begin{document}
\maketitle
\vspace{-2.5em}

\noindent
The tasks are ordered by where the edit goes in the manuscript (front to back).
Tags: \textbf{[R1.4]} = Referee~1's 4th major comment, \textbf{[R2.2]} = Referee~2's
2nd comment, \textbf{[R1.8, minor]} = a minor comment; \emph{(clear/minor)} = quick
edit, \emph{(substantial)} = needs new material. Manuscript numbers/pages are
cross-referenced and update automatically.

\medskip

\noindent\textbf{Progress (2026-06-18): 5 done, 4 partial, 8 open.}
\emph{Label snapshot refreshed 2026-06-18. The new index-wise/tensor-wise
appendix is now Appendix~\ref{app:index_tensor} (the later appendices shift one
letter); the cross-references below already show the new numbering.}

\bigskip

\noindent\textbf{\large Overarching (affects the whole paper)}
\begin{itemize}[leftmargin=1.4em, itemsep=0.45em]
\item \bx \textbf{[R1, summary]} \emph{(substantial)} Make the exposition more
accessible: add a worked example right after each key definition, and an intuitive
proof sketch for the treewidth characterisation (reuse the vertex-elimination /
decomposition-graph slide figure). Clarify the contribution relative to Markov \&
Shi (2008) --- vertex-elimination vs.\ tree-decomposition definition of treewidth,
and that the tensor-wise and index-wise algorithms attain the same optimal
complexity --- and note that the software pre-computes cost under
\texttt{opt-einsum}'s heuristics.
\textbf{(partial done 2026-06-10:} intuitive proof sketch + the new
elimination-sketch figure added after Proposition~\ref{pro:complexity_einsum};
worked examples after the definitions, the Markov--Shi comparison and the
tensor-wise/index-wise point still open.\textbf{)}
\end{itemize}

\medskip

\noindent\textbf{\large By manuscript location (front $\to$ back)}
\begin{enumerate}[label=\arabic*., leftmargin=*, itemsep=0.45em]

\item $\boxtimes$\ \at{p.~\pageref{def:tensor_expression}, after Definitions~\ref{def:tensor_expression}--\ref{def:tensor_contraction} / Example~\ref{eg:running}.}
\emph{[R1.2a, clear]} State that Example~\ref{eg:running} illustrates the case
$B=\emptyset$, and that the examples with $B\neq\emptyset$ are in
Appendix~\ref{app:examples}.
\textbf{(done 2026-06-10}; the added passage also explains why $B \neq
\emptyset$ is indispensable for the proof of
Proposition~\ref{pro:complexity_einsum}; reply written.\textbf{)}

\item $\boxtimes$\ \at{p.~\pageref{eg:running}, Example~\ref{eg:running}.}
\emph{[R1.8, minor]} ``last example'' $\to$ ``running example''.
\textbf{(done; reply written 2026-06-10.)}

\item $\boxtimes$\ \at{p.~\pageref{def:mul_decomp}, Definition~\ref{def:mul_decomp}.}
\emph{[R1.9, minor]} Explain the double-superscript notation $\in[m]^{**}$.
\textbf{(done 2026-06-10:} resolved by \emph{removing} $[m]^{**}$ altogether ---
Definition~\ref{def:mul_decomp} now states $A_k \in [m]^{*}$ directly, matching
Definition~\ref{def:tensor_expression}; reply written.\textbf{)}

\item $\boxtimes$\ \at{Section~\ref{sec:alg_v} (p.~\pageref{sec:alg_v}).}
\emph{[R1.1, substantial]} Add a cost-summary remark (worked on the dcov example)
separating the three sources of cost: tensorization (Algorithm~\ref{alg:tensorization}),
finding the \textsc{Einsum} ordering, and executing it; explain why the entry-access
cost assumed in Definition~\ref{def:tensors} is asymptotically negligible.
\textbf{(done 2026-06-10:} new tensorization-cost remark after
Algorithm~\ref{alg:tensorization} ($O(K n^{m_{\max}})$, dominated since
$m_{\max} \le \mathsf{tw}+1$) + entry-access convention spelled out after
Definition~\ref{def:tensors}; both replies written. The dcov walkthrough from
the original plan was dropped.\textbf{)}

\item \bx \at{Main complexity result, Proposition~\ref{pro:complexity_einsum} (p.~\pageref{pro:complexity_einsum}).}
\emph{[R2.2, substantial]} Add a space/memory-complexity analysis: state it in the
main theorem, add a remark after the tensorization step, and give an appendix bound
on the size of the largest intermediate tensor in the \textsc{Einsum} contraction.
\textbf{(partial done 2026-06-18:} the index-wise space bound $O(n^{\mathsf{tw}})$
is now stated in the main text after Proposition~\ref{pro:complexity_einsum}, and
Appendix~\ref{app:index_tensor} bounds the largest tensor-wise intermediate by
$O(n^{\mathsf{tw}+1})$ (carving-width point) with the dcov necessity example;
still open: a memory column in the empirical tables.\textbf{)}

\item \bx \at{Proposition~\ref{pro:complexity_einsum} (p.~\pageref{pro:complexity_einsum}).}
\emph{[R1.3, clear]} Add simple worked examples illustrating $\mathsf{tw}(G_{\mathcal{A}})$
and how it is computed from the definitions.
\textbf{(partial done 2026-06-10:} proof-sketch paragraphs + the new
elimination-sketch figure are in the manuscript; the worked example computing
$\mathsf{tw}$ for the chain / $K_4$ / dcov graphs is drafted but not yet
pasted in; reply marked partial.\textbf{)}

\item \bx \at{Treewidth result (p.~\pageref{pro:complexity_einsum}).}
\emph{[R2.1, substantial]} Add the characterisation: an $m$-th order $U$-statistic has
a complete decomposition graph iff the maximum treewidth of its quotient graph equals
$m-1$ (easy to check). Even then, after sparsification a single \textsc{Einsum}
operation still beats naive brute force (motif counting is the example).

\item \bx \at{Remark~\ref{rem:ordering} (Remark~3, p.~\pageref{rem:ordering}) and new Appendix~\ref{app:index_tensor}.}
\emph{[R1.4, substantial]} Distinguish the index-wise (theoretical) and tensor-wise
(practical) regimes and report the ``real'' complexity via the HOIF table; add a
remark/experiment comparing greedy vs.\ exhaustive ordering (time to \emph{find} vs.\
FLOPs to \emph{execute}); and compare with the dynamic-programming approach of Zhang
et al.\ (2025).
\textbf{(partial done 2026-06-18:} after Proposition~\ref{pro:complexity_einsum} the
main text now contrasts the index-wise and tensor-wise schemes and works
Example~\ref{eg:index_tensor} ($\mathcal{A}=((1,2),(2,3),(3,4),(4,1),(1,3))$) index-wise;
new Appendix~\ref{app:index_tensor} defines the tensor-wise primitives
(Definition~\ref{def:tw_primitives}), proves the same time upper bound
(Proposition~\ref{pro:tw_time}), and discusses space. Still open: greedy-vs-exhaustive
timing and the Zhang et al.\ (2025) comparison.\textbf{)}

\item \bx \at{Definition~\ref{def:v_set} (Definition~11, p.~\pageref{def:v_set}).}
\emph{[R1.5a, clear]} Add a concrete example of the partition $\pi$ and the meaning of
$Q\in\pi$.

\item \bx \at{Algorithm~\ref{alg:induced_expression} (Algorithm~3, p.~\pageref{alg:induced_expression}).}
\emph{[R1.5b, clear]} Add a small worked example constructing $\mathcal{A}_\pi$ step by
step.

\item \bx \at{Remark~\ref{rem:sparsification_terms_revised} (Remark~4, p.~\pageref{rem:sparsification_terms_revised}).}
\emph{[R1.5c, clear]} Explain how to read equation~\eqref{eq:quotient_set} and how the
stated conclusion follows.

\item \bx \at{Remark~\ref{rem:U-complexity} (Remark~5, p.~\pageref{rem:U-complexity}).}
\emph{[R1.6, substantial]} Explain how $N_l$ and $M$ are computed in practice
(approximate, since exact treewidth is hard) --- e.g.\ extend the HOIF table and give a
complete calculation for the dcov case.

\item \bx \at{Section~\ref{sec:example_hoif} (HOIF, p.~\pageref{sec:example_hoif}).}
\emph{[R1.7a, clear]} Map the HOIF example to the earlier notation: state what
$\mathcal{A}=\{A_1,\dots,A_k\}$ and $h_k(\cdot)$ of Definition~\ref{def:mul_decomp} are
in this case.

\item \bx \at{Observation~\ref{obs:HOIF} (Observation~2, Appendix~\ref{app:HOIF}, p.~\pageref{obs:HOIF}).}
\emph{[R1.7b, substantial]} Give the intuition behind it: why the analysis stops at
$m=12$, and whether an approximate upper bound $M(m)$ captures how the complexity
scales with $m$.

\item \bx \at{Section~\ref{sec:example_motif} (motif counting, p.~\pageref{sec:example_motif}).}
\emph{[R2.3, substantial]} Add a remark: because the relevant decomposition graph is
already complete, our cost depends essentially only on $n$ and is largely independent
of the edge density $p$, whereas enumeration-based methods (igraph, Peregrine,
cuGraph) exploit sparsity and win in the sparse regime.

\item $\boxtimes$\ \at{Appendix~\ref{app:examples} (Appendix~A, p.~\pageref{app:examples}).}
\emph{[R1.2b, clear]} Rename the appendix examples (Example~\ref{eg:matrix multiplication}
and Example~\ref{eg:tensor}) so they no longer clash with the main-text numbering,
e.g.\ A.2 and A.3.
\textbf{(done 2026-06-10} via \texttt{\textbackslash numberwithin\{example\}\{section\}}
after \texttt{\textbackslash appendix}: appendix examples are now A.1, A.2 and
C.1; reply written.\textbf{)}

\end{enumerate}

\end{document}
